\section{When a valid temporal region is still an unsuitable study design}
\label{exec13:temporal-design-review}

The restricted continuous-lag argument separates mathematical coverage
from decision-useful precision. At the true common scalar lag, invariant
Gaussian initialization and independent innovations justify the projected
variance argument. A union over candidate lags need only retain the true
lag on that event; incorrect lags need not themselves produce independent
residuals. This is a useful restricted result. It does not establish a
confidence procedure for an arbitrary lag matrix, unknown order,
nonstationary hidden histories or a general learned partial ancestral
graph. These differences also determine whether a proposed numerical
study is worth conducting \cite{exec13_temporal_design}.

Consider the displayed Chebyshev construction with nine variance
directions, relative tolerance $t=1/4$ and observational error allocation
$\alpha_O=.02$. For a Gaussian sample variance with $\nu$ residual
degrees of freedom, the relative variance is $2/\nu$. Chebyshev and a
union bound therefore give
\[
 \Pr(\hbox{at least one relative deviation exceeds }t)
 \leq \frac{18}{\nu t^2}.
\]
To make this bound at most $.02$ requires $\nu\geq14{,}400$, or
$14{,}401$ transitions under the stated residual convention. With $240$
transitions and $239$ degrees of freedom, the untruncated expression is
approximately $1.20502$ and gives no useful coverage bound. The specified
whole-line fallback is honest about this limitation. Increasing the number
of Monte Carlo replications cannot make that per-request fallback
decision-useful \cite{exec13_temporal_design}.

A sharper pivot is a concrete development route. At the restricted true
lag, an exact chi-square variance pivot can replace the Chebyshev band
while preserving the continuous union. Ordinary development quantiles
at $239$ degrees of freedom and the nine-way $.02$ allocation give
variance multipliers approximately $[.76731,1.34533]$. These values
indicate potential improvement; they are not certified release endpoints.
The implementation must enclose the tail quantiles, observable inputs
and propagated covariance constraints, preserve conservative fallback
under timeout, and profile the continuous-lag calculation. Sampling a
lag grid or accepting a point optimizer's failure does not certify the
whole continuum \cite{exec13_temporal_design}.

The preserved finite-lag study also requires a precise reading. Its
persistent cell returned $507$ finite sets among $512$ requests, so
$507/512=.990234375$ passes the original raw finite-fraction threshold
of $.99$. The five whole-line answers remain visible, but do not justify
declaring that particular threshold failed. Weak relevance and the
unrestricted latent observational cell each returned zero finite sets
and still fail it. False exclusion retained the true graph in zero
requests, despite wide effect sets covering the numerical effect.
These findings concern different obligations: availability, effect
coverage, and scientific validity of the graph class
\cite{exec13_temporal_design}.

A lower confidence-bound rule for availability would be a different
criterion. At family error $.05$ across sixteen comparisons, the
one-sided lower bound for $507/512$ is about $.97130$; even $512/512$
gives about $.98880$. The historical study used a lower-bound screen
for coverage, but a raw fraction for finite availability. They should
not be interchanged retrospectively. A prospective successor may
justify different criteria, provided it states them before final data
and keeps the predecessor's outcomes under their original definitions
\cite{exec13_temporal_design}.

The coverage screen has its own design-power issue. With $512$
independent final requests per cell and error $.05/16$, at least $500$
successes are required for the exact one-sided lower bound to exceed
or equal $.95$. Exact binomial-tail calculations give per-cell passing
probabilities approximately $.001849$, $.234929$, $.770054$ and
$.997670$ when true coverage is respectively $.95$, $.97$, $.98$ and
$.99$. Thus a nominal $.95$ method is very unlikely to certify a
lower bound of $.95$ using this screen. More replications improve
discrimination above the boundary; they do not create high power at
the boundary while retaining the same error control. No product of
these per-cell powers is a joint power calculation without an
across-cell dependence argument \cite{exec13_temporal_design}.

Useful precision also requires an information distinction. In the
unrestricted equal-law latent family, every real effect can produce
the same observed temporal law. No numerical improvement can create
a uniformly finite informative effect set from those unchanged
observations. A direct assigned-action experiment changes the
information available. The proposed Gaussian episode argument assumes
independent episodes, a common transported effect and equal arm
variances. Adjacent records from one trajectory or encouragement in
place of the action require different assumptions and derivations
\cite{exec13_temporal_design}.

The supplied-summary example makes the planning distinction concrete.
At $200$ episodes per arm, reported effect $20$ dollars, reported
pooled variance $100$ dollars squared and the specified conservative
bound, the half-width is about $5.78806$ dollars. For $3{,}000$
eligible customers and an incremental cost of $15$ dollars per offer,
the net-benefit lower endpoint is approximately $-2{,}364.19$ dollars.
Holding the same reported variance and effect fixed, $268$ episodes
per arm first reduce the half-width below $5$ dollars, while $400$
give about $4.08761$. In a future experiment both the estimated effect
and pooled variance would vary. These are conditional arithmetic
screens, not prospective power or recommended bank sample sizes
\cite{exec13_temporal_design}.

The next restricted study therefore needs a certified executable and
input-enclosure policy, measured resource limits, a fixed final grid
of observational and added-information regimes, and an explicit
useful-width or loss criterion. A generated experiment can develop
these elements now. It should preserve strong, weak and zero relevance,
false exclusion, latent equal-law cases and all interval topologies.
Final seeds, replication budgets and stopping rules must follow the
design review. General hidden-state, non-Gaussian, selection and
heterogeneous-transport questions remain unresolved beyond the
restricted claim; public observations alone cannot decide among
observationally equivalent causal worlds
\cite{exec13_temporal_design}.
\label{exec13:temporal-design-review-end}
